% !TeX root = ../抽象代数笔记.tex
% 由主文档通过 \input 载入；请编译主文档。
\setcounter{chapter}{12} % 保留原书章号。
\chapter{多项式环与不可约性}\label{ch:polynomials}

\begin{description}
  \item[学习目标] 掌握域上多项式的除法、分解和不可约性判别，为添加方程的根建立代数基础。
  \item[先修知识] 第\ref{ch:factorization}章，初等多项式运算。
\end{description}

\section{形式多项式、次数与求值}
对交换环 $R$ 定义形式多项式环 $R[x]$，说明形式多项式与其诱导的多项式函数不同；在有限域上不同多项式可能给出同一函数。对整环证明非零多项式乘积的次数可加，并构造求值同态。定义有理函数域 $F(t)=\operatorname{Frac}(F[t])$，其中 $F$ 为域、$t$ 为不定元。

\section{域上多项式的带余除法}
设 $F$ 为域，证明对 $f,g\in F[x]$ 且 $g\ne0$，存在唯一 $q,r\in F[x]$ 使 $f=qg+r$，其中 $r=0$ 或 $\deg r<\deg g$。由此证明 $F[x]$ 为欧几里得整环，建立多项式欧几里得算法、贝祖等式、余式定理与因式定理。

\section{高斯引理与不可约性判别}
在唯一分解整环 $R$ 中定义多项式的内容与本原多项式（primitive polynomial），证明高斯引理（Gauss's lemma），进而证明 $R[x]$ 仍为唯一分解整环。对于正次数的本原多项式，比较它在 $R[x]$ 与 $\operatorname{Frac}(R)[x]$ 中的不可约性。系统安排有理根判别、约化模素数法和艾森斯坦判别法（Eisenstein's criterion），完整说明首项系数与常数项的限制。

\section{根、重根与商环构造域}
证明域上非零 $d$ 次多项式至多有 $d$ 个不同根，定义形式导数与根的重数。对非恒定多项式 $f\in F[x]$，证明它在分裂域中无重根当且仅当 $\gcd(f,f')=1$；分裂域的构造在第\ref{ch:splitting}章完成。证明 $F[x]/(f)$ 为域当且仅当 $f$ 不可约，并说明商环中 $x$ 的剩余类是 $f$ 的一个根。

\section{例题、练习与自检}
判断 $x^4+1$ 在 $\mathbb Q[x]$ 中的不可约性，构造 $\mathbb F_2[x]/(x^2+x+1)$ 并计算元素的逆。比较同一多项式在 $\mathbb Q$、$\mathbb R$ 与 $\mathbb C$ 上的分解。练习中特别安排“没有根但仍可约”的四次多项式，说明无根判别只直接适用于二次、三次多项式。

